% ==============================================================================
% 3. FASILITAS PERUMAHAN, PENDIDIKAN, DAN EKONOMI
% ==============================================================================

\section{Fasilitas Perumahan, Pendidikan, Kesehatan, dan Ekonomi}

Ketersediaan fasilitas infrastruktur dasar, sarana pendidikan, fasilitas pelayanan kesehatan, serta pusat aktivitas ekonomi mikro di Desa Popontolen pada tahun 2025 telah mencerminkan pemerataan pembangunan desa yang berkelanjutan:

\subsection{Akses Kelistrikan, Pendidikan, dan Kesehatan}

\begin{itemize}
  \item \textbf{Kelistrikan:} Sebanyak \textbf{422 keluarga (99,5\%)} telah terlayani aliran listrik PLN mandiri. Hanya terdapat 2 keluarga di Jaga 6 yang belum tersambung listrik PLN dan menjadi sasaran prioritas bantuan desa.
  \item \textbf{Fasilitas Pendidikan:} Layanan pendidikan dasar warga dilayani oleh \textbf{3 unit TK Swasta} dan \textbf{3 unit SD Negeri} yang tersebar di wilayah Jaga 1, Jaga 5, dan Jaga 6.
  \item \textbf{Fasilitas Kesehatan:} Pelayanan kesehatan masyarakat didukung oleh \textbf{1 unit Puskesmas Pembantu (Pustu)} yang berlokasi strategis di Jaga 6 untuk melayani pemeriksaan kesehatan dasar dan posyandu.
  \item \textbf{Telekomunikasi \& Usaha Mikro:} Terdapat \textbf{1 unit Menara BTS} di Jaga 1 yang menopang sinyal 4G/LTE, 1 unit rumah makan, 10 unit warung/kedai makanan-minuman, serta \textbf{18 unit toko/warung kelontong}.
\end{itemize}

\vspace{0.3cm}

\begin{table}[htbp]
\centering
\caption{Distribusi Fasilitas Infrastruktur, Pendidikan, dan Ekonomi (2025)}
\label{tab:fasilitas_desa}
\rowcolors{3}{bpslight}{white}
\begin{tabularx}{\textwidth}{l rrr rrr r}
\toprule
\rowcolor{bpsnavy}
\textcolor{white}{\textbf{Wilayah Jaga}} & 
\textcolor{white}{\textbf{Listrik}} & 
\textcolor{white}{\textbf{No Listrik}} & 
\textcolor{white}{\textbf{TK}} & 
\textcolor{white}{\textbf{SD-N}} & 
\textcolor{white}{\textbf{Kedai}} & 
\textcolor{white}{\textbf{Kelontong}} & 
\textcolor{white}{\textbf{BTS}} \\
\midrule
Jaga 1 & 34 & 0 & 1 & 1 & 0 & 4 & 1 \\
Jaga 2 & 60 & 0 & 0 & 0 & 1 & 4 & 0 \\
Jaga 3 & 38 & 0 & 0 & 0 & 0 & 6 & 0 \\
Jaga 4 & 65 & 0 & 0 & 0 & 3 & 2 & 0 \\
Jaga 5 & 40 & 0 & 1 & 1 & 4 & 0 & 0 \\
Jaga 6 & 45 & \textcolor{red}{\textbf{2}} & 1 & 1 & 2 & 2 & 0 \\
Jaga 7 & 140 & 0 & 0 & 0 & 0 & 0 & 0 \\
\midrule
\rowcolor{bpsborder!40}
\textbf{TOTAL DESA} & \textbf{422} & \textcolor{red}{\textbf{2}} & \textbf{3} & \textbf{3} & \textbf{10} & \textbf{18} & \textbf{1} \\
\bottomrule
\end{tabularx}
\end{table}

\subsection{Profil Inklusivitas dan Penyandang Disabilitas}

Sebagai bentuk komitmen terhadap tata kelola desa yang inklusif dan ramah kelompok rentan, pendataan potensi desa mencatat keberadaan \textbf{24 orang penyandang disabilitas} di Desa Popontolen pada tahun 2025 yang tersebar di 7 wilayah Jaga.

\vspace{0.3cm}

\begin{table}[htbp]
\centering
\caption{Sebaran Penyandang Disabilitas Menurut Ragam Disabilitas per Jaga (2025)}
\label{tab:disabilitas}
\rowcolors{3}{bpslight}{white}
\begin{tabularx}{\textwidth}{l rrrrr r}
\toprule
\rowcolor{bpsnavy}
\textcolor{white}{\textbf{Wilayah Jaga}} & 
\textcolor{white}{\textbf{Fisik}} & 
\textcolor{white}{\textbf{Mental}} & 
\textcolor{white}{\textbf{Netra}} & 
\textcolor{white}{\textbf{Wicara}} & 
\textcolor{white}{\textbf{Rungu}} & 
\textcolor{white}{\textbf{Total Jiwa}} \\
\midrule
Jaga 1 & 1 & 0 & 0 & 0 & 0 & \textbf{1} \\
Jaga 2 & 2 & 1 & 1 & 0 & 0 & \textbf{4} \\
Jaga 3 & 1 & 1 & 0 & 0 & 0 & \textbf{2} \\
Jaga 4 & 0 & 2 & 0 & 1 & 0 & \textbf{3} \\
Jaga 5 & 0 & 1 & 0 & 2 & 1 & \textbf{4} \\
Jaga 6 & 2 & 0 & 0 & 0 & 0 & \textbf{2} \\
Jaga 7 & 1 & 4 & 2 & 0 & 0 & \textbf{7} \\
\midrule
\rowcolor{bpsborder!40}
\textbf{TOTAL DESA} & \textbf{8} & \textbf{9} & \textbf{3} & \textbf{3} & \textbf{1} & \textbf{24} \\
\bottomrule
\end{tabularx}
\end{table}

Data disabilitas ini menjadi landasan bagi Pemerintah Desa Popontolen dalam menyusun program jaminan sosial, bantuan alat bantu mobilitas, serta program inklusi pemberdayaan masyarakat rentan yang terintegrasi dalam APBDes.
